This phase turns the basin-entry story into a separatrix-margin test. Using the Phase~B PG basin mask, each initial policy receives a signed distance to the nearest opposite PG-basin cell. For PG-fail cells this is the empirical distance that an early update must overcome before ordinary PG should be expected to finish the cooperative selection.

\begin{figure}[h]
  \centering
  \includegraphics[width=0.95\linewidth]{phase_w_separatrix_margin}
  \caption{Separatrix-margin certificate. Left: signed distance to the PG basin boundary. Middle: first-step peer displacement versus boundary distance. Right: gain probability after normalizing the peer projection by the empirical boundary distance.}
  \label{fig:phase-w-margin}
\end{figure}

\begin{table}[h]
  \centering
\begin{tabular}{lcc}
  \toprule
  Metric & Peer / margin & Improvement / margin \\
  \midrule
  AUC on PG-fail cells & 0.954 & 0.959 \\
  Spearman with gain label & 0.654 & 0.660 \\
\bottomrule
\end{tabular}
  \caption{Margin-normalized prediction of Meta-MAPG gains.}
  \label{tab:phase-w-margin}
\end{table}

The median boundary distance is 0.051 for gained PG-fail cells and 0.308 for non-gained PG-fail cells. This is stricter than Phase~R: the peer term must be large in the right direction relative to the local separatrix margin, not merely positive.
